\documentclass{zhvt-classic}

% 使用本地源流明体（GenRyuMin2 TW）按文件路径加载，不依赖系统字体名
% 版框（内框／界栏／书口侧栏／鱼尾）改用与版心印章同一印泥棕色：
%   印章图片 yshs.png 实色 = RGB(135,68,52) = HTML 874434
\usepackage{xcolor}
\usepackage{tikz}
% sealbrown 双注册：xcolor 供 TikZ/pgf，L3 供 cls 版框
\definecolor{sealbrown}{HTML}{874434}
\ExplSyntaxOn
\color_set:nnn { sealbrown } { HTML } { 874434 }
\ExplSyntaxOff

\zhvtset{
  grid color       = sealbrown,
  main font        = GenRyuMin2TW-SB,
  main font path   = C:/Users/Administrator/texmf/fonts/opentype/public/GenRyuMin2TW/,
jiazhu font      = GenRyuMin2TW-M,
  jiazhu font path = C:/Users/Administrator/texmf/fonts/opentype/public/GenRyuMin2TW/,
  % 页码框容纳 4 字：正常页「九十七」+前置空 1 格=4 行恰好；
  % 100+ 页「一百零一」共 5 行超出，末字「一/二/三」被框裁掉（2026-10-04 实测），故放宽到 5
  shukou page number chars = 5,
  % 书口书名框高从默认 6 字减到 4 字，把省出的 2 字高拨给 chapter 框，
  % 使 6 字书名+卷名放进 9 字高的 chapter 框后不再压线（2026-10-05）
  shukou title chars = 4,
}

\title{史记}[魏世家]
\author{汉·司马迁}
\maker{闲叟漫记·重製}


% ===== 索隐述赞样式（2026-11-08 四次定稿 · 愚大人亲手改定参数）=====
% ⚠ 竖排前提：自定义字体家族必须带 RawFeature={vertical:+vert}；
%   simhei.ttf 加 +vhal 会整体加载失败静默回退，只写 +vert；simfang 同理。
\newCJKfontfamily{\shuzanfont}[Path=C:/Windows/Fonts/,RawFeature={vertical:+vert}]{simhei.ttf}
\newCJKfontfamily{\shuzanbody}[Path=C:/Windows/Fonts/,RawFeature={vertical:+vert}]{simfang.ttf}
% 题头：152x30pt 圆角 6pt 棕底白字；baseline (0,2) 使整框落入界行格子不压版框线；
% 白字节点 x 不补偿、y +8.2pt（像素级墨迹实测 dx/dy≈0，L/R 留白约 20pt）。
\newcommand{\shuzanhead}[1]{%
  \begin{tikzpicture}[baseline={(0,2)}]
    \fill[sealbrown, rounded corners=6pt]
      (-76pt,-15pt) rectangle (76pt,15pt);
    \node[text=white, inner sep=0pt, anchor=center] at (0pt,8.2pt)
      {\shuzanfont\fontsize{23}{23}\selectfont #1};
  \end{tikzpicture}%
}
% ★ 修正书口卷名位置：chapter 框起始位置改回 cls 默认值，避免全局影响。
%   只在每个卷首页通过 \juan 命令里的空命令微调该页位置（安全，不会跨页污染）。
\ExplSyntaxOn
\AtBeginDocument{
  \dim_set:Nn \l__zhvt_shukou_chapter_title_start_dim
    { \l__zhvt_shukou_title_height_dim + \l__zhvt_font_size_dim * 2 }
  \dim_set:Nn \l__zhvt_shukou_page_number_start_dim
    { \l__zhvt_shukou_chapter_title_start_dim + \l__zhvt_shukou_chapter_title_height_dim }
}
\ExplSyntaxOff

\begin{document}

% ★★ 版心（书口）排布覆写：把书名放到鱼尾之下（2026-10-04）
%
%   cls 默认版心自上而下：书名（title 静态框，在鱼尾之上）→ 鱼尾 → 篇名（chapter 动态框）→ 页码
%   本模板改为：          鱼尾（位置不动）→ 书名 → 篇名 → 页码分割线 → 页码
%
%   手法：把 cls 放在 title 静态框里的书名，改放进 chapter 动态框（鱼尾之下）；
%         title 静态框留空。书名取 \l__zhvt_title_str，即 \title{...} 的第一参数，
%         所以换书时只改 \title，不必在此重复写书名。
%
%   ⚠ 三条必须遵守的写法（2026-10-04 实测）：
%     ① 汉字必须用 \hbox{...} 包住。flowfram 会把框内容装进一个宽 14.4mm（= 一个 line_sep
%        = 书口栏宽）的纵向 minipage 并按该宽度折行，而整页靠 PDF /Rotate 90 显示、
%        框内文字本身不变向 —— 不包 \hbox 时「围炉夜话」4 字会被逐字折成 4 行，
%        在版面上就成了横排 4 字并横向溢出版框。
%     ② 竖排下**一个 \hbox 就是一整列**。书名与篇名要同处一列，必须放进同一个 \hbox；
%        分成两个 \hbox 时，后一个会被当成新的一列、甩到隔壁栏去
%        （实测「序」跑到书口栏外侧 14.4mm 的正文列上）。中间的 \quad 在框内是横向空白，
%        显示出来正好是书名与篇名之间的纵向间隔。
%     ③ ★ 内容盒要写成 **与框等宽的定宽盒** `\hbox to \l__zhvt_shukou_chapter_title_height_dim{…\hfil}`。
%        实测：chapter 框是把内容盒沿纵向**居中**摆放的，内容宽度一变，书名就跟着漂：
%        自然宽度的 \hbox（书名 38.4mm）落在 x=157.2mm，加篇名后（57.6mm）落在 x=152.4mm，
%        目录页（无篇名）与正文页的书名就差半格。做成等宽盒后，无论有没有篇名，
%        内容固定从框顶起排，实测两页均落在 x=142.8mm，不再漂。
\ExplSyntaxOn
% 把书口 chapter 框（书名+篇名）起始位置下移 1 字高，避免 4 字卷名顶到鱼尾
% ⚠ cls 的 \__zhvt_calculate_dimensions: 在 \AtBeginDocument 里执行，会重算该 dim；
%   必须把覆写挂到 begin document 之后，否则会被二次覆盖（2026-10-05 实测）。
% (旧方案已移到导言区，直接修改 chapter 框起始位置)
% 卷标题命令 #1=书口篇名(鱼尾) #2=正文卷名（含卷序，卷名内空格用 ~）
%
% ★★ 目录页两项修订（2026-10-05，愚大人指出）★★
%   ① 页码起点统一对齐：
%      原状 —— 目录项写成 \par\hspace{4em}卷名\hspace{4em}页码，页码起点 = 卷名实宽 + 4em。
%      卷名实宽随字数变：「卷一~梁惠王上」6 字 +~、「卷七~离娄上」5 字 +~，
%      差 1 个字位 → 卷七~卷十 4 条的页码整体比其余 10 条高 27.21pt（一个字），看着上下不齐。
%      改法 —— 卷名补全角空格（宽 = 1 字 = 27.215pt）补足到 **7 字符位**再写进目录，
%      14 条页码起点遂全部落在同一条竖线上（起点对齐；末字随页码字数自然参差）。
%      ⚠ 只对目录里的副本 \l_tmpc_tl 补，正文卷名仍用 \l_tmpa_tl 原样 —— 正文像素不受影响。
%   ② 卷名整体上移半字位：前导 \hspace 由 4em 改 3.5em。
%      实测原状目录块在版框内偏下（天头留白 139.6pt / 地脚留白 110.7pt），
%      上移 0.5em = 13.6pt 后上下基本均分，视觉居中。
\NewDocumentCommand { \juan } { mm } {
	\clearpage
	\tl_set:Nn \l__zhvt_chapter_title_tl { #1 }
	\tl_set:Nn \l_tmpa_tl { #2 }
	\tl_set_eq:NN \l_tmpc_tl \l_tmpa_tl
	\int_step_inline:nn { 7 - \str_count:n { #2 } }
	  { \tl_put_right:Nn \l_tmpc_tl { 　 } }
	\tl_set:Nx \l_tmpb_tl { \zhnum { page } }
	\tl_put_left:Nn \l_tmpb_tl { \hspace { 4em } }
	\tl_put_left:Nn \l_tmpb_tl { \l_tmpc_tl }
	\tl_put_left:Nn \l_tmpb_tl { \par \hspace { 3.5em } }
	\addtocontents { toc } { \l_tmpb_tl }
	\skip_horizontal:n { \l__zhvt_font_size_dim * 4 }
	\l_tmpa_tl
	\par
	\thispagestyle { empty }
}
\cs_set_protected:Npn \__zhvt_set_shukou_content: {
	% ★ 鱼尾上方留白（title 静态框）嵌入闲叟钤印：
	%   整页 /Rotate 90，故此处把图用 \rotatebox{90} 转正为竖章；
	%   尺寸卡在版心条幅内（书内约 7mm 宽 × 33.8mm 高，留余量不压鱼尾）。
	\setstaticcontents * { title } {
		\vbox_to_ht:nn { \l__zhvt_baseline_skip_dim } {
			\vskip 14.2pt
			\hbox_to_wd:nn { \l__zhvt_shukou_title_height_dim } {
				\hfil
				\rotatebox [ origin = c ] { 90 } { \includegraphics [ width = 7mm , height = 33.8mm ] { yshs.png } }
				\hfil
			}
		}
	}
% ★ 版心（书口）统一改法（2026-11-16 愚大人定）：18pt＋框内居中＋书口文字同版框色
%   ① 原 27.2pt 时 7 字长篇名（如「廉颇蔺相如列传」）末字溢出 chapter 框、压住页码；
%      改 \fontsize{18pt}{18pt}，前后 \hfil 令字串在框内居中（不随篇名字数漂移）。
%   ② \vbox{\vskip 8.2pt} 定量下移，抵消缩字号后栏宽方向 8.2pt 上偏（像素实测 y 中心 543.3 = 原版）。
%      ⚠ \raisebox / \vspace* 在 flowfram 框内无效；换字号须同步换补偿：16pt→10.1、19pt→7.4。
%   ③ chapter 与 page~number 均置 \color{sealbrown}（HTML 874434），与版框／界行同色；正文仍黑。
%   ⚠ 页码字位、钤印、鱼尾、界行、正文一律未动；五部合订本同此口径。
	\setdynamiccontents * { chapter }
	{ \color{sealbrown} \vbox { \vskip 8.2pt \hbox to \l__zhvt_shukou_chapter_title_height_dim
		{ \hfil \fontsize{18pt}{18pt}\selectfont \l__zhvt_title_str \quad \l__zhvt_chapter_title_tl \hfil } } }
	\setdynamiccontents * { page~number } { \color{sealbrown} \quad \zhnum { page } }
}
\__zhvt_set_shukou_content:
\ExplSyntaxOff

\maketitle

% 单篇，无多卷目录，不渲染目录页

\juan{魏世家}{卷四十四~魏世家第十四}
\hspace*{2em}魏之先。毕公高之後也。毕公高与周同姓。武王之伐纣。而高封於毕。於是为毕姓。其後绝封。为庶人。或在中国。或在夷狄。其苗裔曰毕万。事晋献公。\\
\hspace*{2em}献公之十六年。赵夙为御。毕万为右。以伐霍。耿。魏。灭之。以耿封赵夙。以魏封毕万。为大夫。卜偃曰。﹁毕万之後必大矣。万。满数也。魏。大名也。以是始赏。天开之矣。天子曰兆民。诸侯曰万民。今命之大。以从满数。其必有众。﹂初。毕万卜事晋。遇屯之比。辛廖占之。曰。﹁吉。屯固比入。吉孰大焉。其必蕃昌。﹂\\
\hspace*{2em}毕万封十一年。晋献公卒。四子争更立。晋乱。而毕万之世弥大。从其国名为魏氏。生武子。魏武子以魏诸子事晋公子重耳。晋献公之二十一年。武子从重耳出亡。十九年反。重耳立为晋文公。而令魏武子袭魏氏之後封。列为大夫。治於魏。生悼子。\\
\hspace*{2em}魏悼子徙治霍。生魏绛。\\
\hspace*{2em}魏绛事晋悼公。悼公三年。会诸侯。悼公弟杨干乱行。魏绛僇辱杨干。悼公怒曰。﹁合诸侯以为荣。今辱吾弟。﹂将诛魏绛。或说悼公。悼公止。卒任魏绛政。使和戎。翟。戎。翟亲附。悼公之十一年。曰。﹁自吾用魏绛。八年之中。九合诸侯。戎。翟和。子之力也。﹂赐之乐。三让。然後受之。徙治安邑。魏绛卒。谥为昭子。生魏嬴。嬴生魏献子。\\
\hspace*{2em}献子事晋昭公。昭公卒而六卿彊。公室卑。\\
\hspace*{2em}晋顷公之十二年。韩宣子老。魏献子为国政。晋宗室祁氏。羊舌氏相恶。六卿诛之。尽取其邑为十县。六卿各令其子为之大夫。献子与赵简子。中行文子。范献子并为晋卿。\\
\hspace*{2em}其後十四岁而孔子相鲁。後四岁。赵简子以晋阳之乱也。而与韩。魏共攻范。中行氏。魏献子生魏侈。魏侈与赵鞅共攻范。中行氏。\\
\hspace*{2em}魏侈之孙曰魏桓子。与韩康子。赵襄子共伐灭知伯。分其地。\\
\hspace*{2em}桓子之孙曰文侯都。魏文侯元年。秦灵公之元年也。与韩武子。赵桓子。周威王同时。\\
\hspace*{2em}六年。城少梁。十三年。使子击围繁。庞。出其民。十六年。伐秦。筑临晋元里。\\
\hspace*{2em}十七年。伐中山。使子击守之。赵仓唐傅之。子击逢文侯之师田子方於朝歌。引车避。下谒。田子方不为礼。子击因问曰。﹁富贵者骄人乎。且贫贱者骄人乎。﹂子方曰。﹁亦贫贱者骄人耳。夫诸侯而骄人则失其国。大夫而骄人则失其家。贫贱者。行不合。言不用。则去之楚。越。若脱鵕然。柰何其同之哉。﹂子击不怿而去。西攻秦。至郑而还。筑雒阴。合阳。\\
\hspace*{2em}二十二年。魏。赵。韩列为诸侯。\\
\hspace*{2em}二十四年。秦伐我。至阳狐。\\
\hspace*{2em}二十五年。子击生子。\\
\hspace*{2em}文侯受子夏经艺。客段干木。过其闾。未尝不轼也。秦尝欲伐魏。或曰。﹁魏君贤人是礼。国人称仁。上下和合。未可图也。﹂文侯由此得誉於诸侯。\\
\hspace*{2em}任西门豹守鄴。而河内称治。\\
\hspace*{2em}魏文侯谓李克曰。﹁先生尝教寡人曰﹃家贫则思良妻。国乱则思良相﹄。今所置非成则璜。二子何如。﹂李克对曰。﹁臣闻之。卑不谋尊。疏不谋戚。臣在阙门之外。不敢当命。﹂文侯曰。﹁先生临事勿让。﹂李克曰。﹁君不察故也。居视其所亲。富视其所与。达视其所举。穷视其所不为。贫视其所不取。五者足以定之矣。何待克哉。﹂文侯曰。﹁先生就舍。寡人之相定矣。﹂李克趋而出。过翟璜之家。翟璜曰。﹁今者闻君召先生而卜相。果谁为之。﹂李克曰。﹁魏成子为相矣。﹂翟璜忿然作色曰。﹁以耳目之所睹记。臣何负於魏成子。西河之守。臣之所进也。君内以鄴为忧。臣进西门豹。君谋欲伐中山。臣进乐羊。中山以拔。无使守之。臣进先生。君之子无傅。臣进屈侯鲋。臣何以负於魏成子。﹂李克曰。﹁且子之言克於子之君者。岂将比周以求大官哉。君问而置相﹃非成则璜。二子何如﹄。克对曰。﹃君不察故也。居视其所亲。富视其所与。达视其所举。穷视其所不为。贫视其所不取。五者足以定之矣。何待克哉。﹄是以知魏成子之为相也。且子安得与魏成子比乎。魏成子以食禄千锺。什九在外。什一在内。是以东得卜子夏。田子方。段干木。此三人者。君皆师之。子之所进五人者。君皆臣之。子恶得与魏成子比也。﹂翟璜逡巡再拜曰。﹁璜。鄙人也。失对。原卒为弟子。﹂\\
\hspace*{2em}二十六年。虢山崩。壅河。\\
\hspace*{2em}三十二年。伐郑。城酸枣。败秦于注。三十五年。齐伐取我襄陵。三十六年。秦侵我阴晋。\\
\hspace*{2em}三十八年。伐秦。败我武下。得其将识。是岁。文侯卒。子击立。是为武侯。\\
\hspace*{2em}魏武侯元年。赵敬侯初立。公子朔为乱。不胜。奔魏。与魏袭邯郸。魏败而去。\\
\hspace*{2em}二年。城安邑。王垣。\\
\hspace*{2em}七年。伐齐。至桑丘。九年。翟败我于澮。使吴起伐齐。至灵丘。齐威王初立。\\
\hspace*{2em}十一年。与韩。赵三分晋地。灭其後。\\
\hspace*{2em}十三年。秦献公县栎阳。十五年。败赵北蔺。\\
\hspace*{2em}十六年。伐楚。取鲁阳。武侯卒。子立。是为惠王。\\
\hspace*{2em}惠王元年。初。武侯卒也。子与公中缓争为太子。公孙颀自宋入赵。自赵入韩。谓韩懿侯曰。﹁魏与公中缓争为太子。君亦闻之乎。今魏得王错。挟上党。固半国也。因而除之。破魏必矣。不可失也。﹂懿侯说。乃与赵成侯合军并兵以伐魏。战于浊泽。魏氏大败。魏君围。赵谓韩曰。﹁除魏君。立公中缓。割地而退。我且利。﹂韩曰。﹁不可。杀魏君。人必曰暴。割地而退。人必曰贪。不如两分之。魏分为两。不彊於宋。卫。则我终无魏之患矣。﹂赵不听。韩不说。以其少卒夜去。惠王之所以身不死。国不分者。二家谋不和也。若从一家之谋。则魏必分矣。故曰﹁君终无適子。其国可破也﹂。\\
\hspace*{2em}二年。魏败韩于马陵。败赵于怀。三年。齐败我观。五年。与韩会宅阳。城武堵。为秦所败。六年。伐取宋仪台。九年。伐败韩于澮。与秦战少梁。虏我将公孙痤。取庞。秦献公卒。子孝公立。\\
\hspace*{2em}十年。伐取赵皮牢。彗星见。十二年。星昼坠。有声。\\
\hspace*{2em}十四年。与赵会鄗。十五年。鲁。卫。宋。郑君来朝。十六年。与秦孝公会杜平。侵宋黄池。宋复取之。\\
\hspace*{2em}十七年。与秦战元里。秦取我少梁。围赵邯郸。十八年。拔邯郸。赵请救于齐。齐使田忌。孙膑救赵。败魏桂陵。\\
\hspace*{2em}十九年。诸侯围我襄陵。筑长城。塞固阳。\\
\hspace*{2em}二十年。归赵邯郸。与盟漳水上。二十一年。与秦会彤。赵成侯卒。二十八年。齐威王卒。中山君相魏。\\
\hspace*{2em}三十年。魏伐赵。赵告急齐。齐宣王用孙子计。救赵击魏。魏遂大兴师。使庞涓将。而令太子申为上将军。过外黄。外黄徐子谓太子曰。﹁臣有百战百胜之术。﹂太子曰。﹁可得闻乎。﹂客曰。﹁固原效之。﹂曰。﹁太子自将攻齐。大胜并莒。则富不过有魏。贵不益为王。若战不胜齐。则万世无魏矣。此臣之百战百胜之术也。﹂太子曰。﹁诺。请必从公之言而还矣。﹂客曰。﹁太子虽欲还。不得矣。彼劝太子战攻。欲啜汁者众。太子虽欲还。恐不得矣。﹂太子因欲还。其御曰。﹁将出而还。与北同。﹂太子果与齐人战。败於马陵。齐虏魏太子申。杀将军涓。军遂大破。\\
\hspace*{2em}三十一年。秦。赵。齐共伐我。秦将商君诈我将军公子卬而袭夺其军。破之。秦用商君。东地至河。而齐。赵数破我。安邑近秦。於是徙治大梁。以公子赫为太子。\\
\hspace*{2em}三十三年。秦孝公卒。商君亡秦归魏。魏怒。不入。三十五年。与齐宣王会平阿南。\\
\hspace*{2em}惠王数被於军旅。卑礼厚币以招贤者。邹衍。淳于髡。孟轲皆至梁。梁惠王曰。﹁寡人不佞。兵三折於外。太子虏。上将死。国以空虚。以羞先君宗庙社稷。寡人甚丑之。叟不远千里。辱幸至弊邑之廷。将何利吾国。﹂孟轲曰。﹁君不可以言利若是。夫君欲利则大夫欲利。大夫欲利则庶人欲利。上下争利。国则危矣。为人君。仁义而已矣。何以利为。﹂\\
\hspace*{2em}三十六年。复与齐王会甄。是岁。惠王卒。子襄王立。\\
\hspace*{2em}襄王元年。与诸侯会徐州。相王也。追尊父惠王为王。\\
\hspace*{2em}五年。秦败我龙贾军四万五千于雕阴。围我焦。曲沃。予秦河西之地。\\
\hspace*{2em}六年。与秦会应。秦取我汾阴。皮氏。焦。魏伐楚。败之陉山。七年。魏尽入上郡于秦。秦降我蒲阳。八年。秦归我焦。曲沃。\\
\hspace*{2em}十二年。楚败我襄陵。诸侯执政与秦相张仪会齧桑。十三年。张仪相魏。魏有女子化为丈夫。秦取我曲沃。平周。\\
\hspace*{2em}十六年。襄王卒。子哀王立。张仪复归秦。\\
\hspace*{2em}哀王元年。五国共攻秦。不胜而去。\\
\hspace*{2em}二年。齐败我观津。五年。秦使樗里子伐取我曲沃。走犀首岸门。六年。秦来立公子政为太子。与秦会临晋。七年。攻齐。与秦伐燕。\\
\hspace*{2em}八年。伐卫。拔列城二。见卫君曰。﹁请罢魏兵。免成陵君可乎。﹂卫君曰。﹁先生果能。孤请世世以卫事先生。﹂如耳见成陵君曰。﹁昔者魏伐赵。断羊肠。拔阏与。约斩赵。赵分而为二。所以不亡者。魏为从主也。今卫已迫亡。将西请事於秦。与其以秦醳卫。不如以魏醳卫。卫之德魏必终无穷。﹂成陵君曰。﹁诺。﹂如耳见魏王曰。﹁臣有谒於卫。卫故周室之别也。其称小国。多宝器。今国迫於难而宝器不出者。其心以为攻卫醳卫不以王为主。故宝器虽出必不入於王也。臣窃料之。先言醳卫者必受卫者也。﹂如耳出。成陵君入。以其言见魏王。魏王听其说。罢其兵。免成陵君。终身不见。\\
\hspace*{2em}九年。与秦王会临晋。张仪。魏章皆归于魏。魏相田需死。楚害张仪。犀首。薛公。楚相昭鱼谓苏代曰。﹁田需死。吾恐张仪。犀首。薛公有一人相魏者也。﹂代曰。﹁然相者欲谁而君便之。﹂昭鱼曰。﹁吾欲太子之自相也。﹂代曰。﹁请为君北。必相之。﹂昭鱼曰。﹁柰何。﹂对曰。﹁君其为梁王。代请说君。﹂昭鱼曰。﹁柰何。﹂对曰。﹁代也从楚来。昭鱼甚忧。曰。﹃田需死。吾恐张仪。犀首。薛公有一人相魏者也。﹄代曰。﹃梁王。长主也。必不相张仪。张仪相。必右秦而左魏。犀首相。必右韩而左魏。薛公相。必右齐而左魏。梁王。长主也。必不便也。﹄王曰。﹃然则寡人孰相。﹄代曰。﹃莫若太子之自相。太子之自相。是三人者皆以太子为非常相也。皆将务以其国事魏。欲得丞相玺也。以魏之彊。而三万乘之国辅之。魏必安矣。故曰莫若太子之自相也。﹄﹂遂北见梁王。以此告之。太子果相魏。\\
\hspace*{2em}十年。张仪死。十一年。与秦武王会应。十二年。太子朝於秦。秦来伐我皮氏。未拔而解。十四年。秦来归武王后。十六年。秦拔我蒲反。阳晋。封陵。十七年。与秦会临晋。秦予我蒲反。十八年。与秦伐楚。`二十一年。与齐。韩共败秦军函谷。\\
\hspace*{2em}二十三年。秦复予我河外及封陵为和。哀王卒。子昭王立。\\
\hspace*{2em}昭王元年。秦拔我襄城。二年。与秦战。我不利。三年。佐韩攻秦。秦将白起败我军伊阙二十四万。六年。予秦河东地方四百里。芒卯以诈重。七年。秦拔我城大小六十一。八年。秦昭王为西帝。齐湣王为东帝。月馀。皆复称王归帝。九年。秦拔我新垣。曲阳之城。\\
\hspace*{2em}十年。齐灭宋。宋王死我温。十二年。与秦。赵。韩。燕共伐齐。败之济西。湣王出亡。燕独入临菑。与秦王会西周。\\
\hspace*{2em}十三年。秦拔我安城。兵到大梁。去。十八年。秦拔郢。楚王徙陈。\\
\hspace*{2em}十九年。昭王卒。子安釐王立。\\
\hspace*{2em}安釐王元年。秦拔我两城。二年。又拔我二城。军大梁下。韩来救。予秦温以和。三年。秦拔我四城。斩首四万。四年。秦破我及韩。赵。杀十五万人。走我将芒卯。魏将段干子请予秦南阳以和。苏代谓魏王曰。﹁欲玺者段干子也。欲地者秦也。今王使欲地者制玺。使欲玺者制地。魏氏地不尽则不知已。且夫以地事秦。譬犹抱薪救火。薪不尽。火不灭。﹂王曰。﹁是则然也。虽然。事始已行。不可更矣。﹂对曰。﹁王独不见夫博之所以贵枭者。便则食。不便则止矣。今王曰﹃事始已行。不可更﹄。是何王之用智不如用枭也。﹂\\
\hspace*{2em}九年。秦拔我怀。十年。秦太子外质於魏死。十一年。秦拔我郪丘。\\
\hspace*{2em}秦昭王谓左右曰。﹁今时韩。魏与始孰彊。﹂对曰。﹁不如始彊。﹂王曰。﹁今时如耳。魏齐与孟尝。芒卯孰贤。﹂对曰。﹁不如。﹂王曰。﹁以孟尝。芒卯之贤。率彊韩。魏以攻秦。犹无柰寡人何也。今以无能之如耳。魏齐而率弱韩。魏以伐秦。其无柰寡人何亦明矣。﹂左右皆曰。﹁甚然。﹂中旗冯琴而对曰。﹁王之料天下过矣。当晋六卿之时。知氏最彊。灭范。中行。又率韩。魏之兵以围赵襄子於晋阳。决晋水以灌晋阳之城。不湛者三版。知伯行水。魏桓子御。韩康子为参乘。知伯曰。﹃吾始不知水之可以亡人之国也。乃今知之。﹄汾水可以灌安邑。绛水可以灌平阳。魏桓子肘韩康子。韩康子履魏桓子。肘足接於车上。而知氏地分。身死国亡。为天下笑。今秦兵虽彊。不能过知氏。韩。魏虽弱。尚贤其在晋阳之下也。此方其用肘足之时也。原王之勿易也。﹂於是秦王恐。\\
\hspace*{2em}齐。楚相约而攻魏。魏使人求救於秦。冠盖相望也。而秦救不至。魏人有唐雎者。年九十馀矣。谓魏王曰。﹁老臣请西说秦王。令兵先臣出。﹂魏王再拜。遂约车而遣之。唐雎到。入见秦王。秦王曰。﹁丈人芒然乃远至此。甚苦矣。夫魏之来求救数矣。寡人知魏之急已。﹂唐雎对曰。﹁大王已知魏之急而救不发者。臣窃以为用策之臣无任矣。夫魏。一万乘之国也。然所以西面而事秦。称东籓。受冠带。祠春秋者。以秦之彊足以为与也。今齐。楚之兵已合於魏郊矣。而秦救不发。亦将赖其未急也。使之大急。彼且割地而约从。王尚何救焉。必待其急而救之。是失一东籓之魏而彊二敌之齐。楚。则王何利焉。﹂於是秦昭王遽为发兵救魏。魏氏复定。\\
\hspace*{2em}赵使人谓魏王曰。﹁为我杀范痤。吾请献七十里之地。﹂魏王曰。﹁诺。﹂使吏捕之。围而未杀。痤因上屋骑危。谓使者曰。﹁与其以死痤市。不如以生痤市。有如痤死。赵不予王地。则王将柰何。故不若与先定割地。然後杀痤。﹂魏王曰。﹁善。﹂痤因上书信陵君曰。﹁痤。故魏之免相也。赵以地杀痤而魏王听之。有如彊秦亦将袭赵之欲。则君且柰何。﹂信陵君言於王而出之。\\
\hspace*{2em}魏王以秦救之故。欲亲秦而伐韩。以求故地。无忌谓魏王曰。\\
\hspace*{2em}秦与戎翟同俗。有虎狼之心。贪戾好利无信。不识礼义德行。苟有利焉。不顾亲戚兄弟。若禽兽耳。此天下之所识也。非有所施厚积德也。故太后母也。而以忧死。穰侯舅也。功莫大焉。而竟逐之。两弟无罪。而再夺之国。此於亲戚若此。而况於仇雠之国乎。今王与秦共伐韩而益近秦患。臣甚惑之。而王不识则不明。群臣莫以闻则不忠。\\
\hspace*{2em}今韩氏以一女子奉一弱主。内有大乱。外交彊秦魏之兵。王以为不亡乎。韩亡。秦有郑地。与大梁鄴。王以为安乎。王欲得故地。今负彊秦之亲。王以为利乎。\\
\hspace*{2em}秦非无事之国也。韩亡之後必将更事。更事必就易与利。就易与利必不伐楚与赵矣。是何也。夫越山逾河。绝韩上党而攻彊赵。是复阏与之事。秦必不为也。若道河内。倍鄴。朝歌。绝漳滏水。与赵兵决於邯郸之郊。是知伯之祸也。秦又不敢。伐楚。道涉谷。行三千里。而攻冥戹之塞。所行甚远。所攻甚难。秦又不为也。若道河外。倍大梁。右上蔡。召陵。与楚兵决於陈郊。秦又不敢。故曰秦必不伐楚与赵矣。又不攻卫与齐矣。\\
\hspace*{2em}夫韩亡之後。兵出之日。非魏无攻已。秦固有怀。茅。邢丘。城垝津以临河内。河内共。汲。必危。有郑地。得垣雍。决荧泽水灌大梁。大梁必亡。王之使者出过而恶安陵氏於秦。秦之欲诛之久矣。秦叶阳。昆阳与舞阳邻。听使者之恶之。随安陵氏而亡之。绕舞阳之北。以东临许。南国必危。国无害乎。\\
\hspace*{2em}夫憎韩不爱安陵氏可也。夫不患秦之不爱南国非也。异日者。秦在河西晋。国去梁千里。有河山以阑之。有周韩以间之。从林乡军以至于今。秦七攻魏。五入囿中。边城尽拔。文台堕。垂都焚。林木伐。麋鹿尽。而国继以围。又长驱梁北。东至陶卫之郊。北至平监。所亡於秦者。山南山北。河外河内。大县数十。名都数百。秦乃在河西晋。去梁千里。而祸若是矣。又况於使秦无韩。有郑地。无河山而阑之。无周韩而间之。去大梁百里。祸必由此矣。\\
\hspace*{2em}异日者。从之不成也。楚。魏疑而韩不可得也。今韩受兵三年。秦桡之以讲。识亡不听。投质於赵。请为天下雁行顿刃。楚。赵必集兵。皆识秦之欲无穷也。非尽亡天下之国而臣海内。必不休矣。是故臣原以从事王。王速受楚赵之约。而挟韩之质以存韩。而求故地。韩必效之。\\
\hspace*{2em}夫存韩安魏而利天下。此亦王之天时已。通韩上党於共。甯。使道安成。出入赋之。是魏重质韩以其上党也。今有其赋。足以富国。韩必德魏爱魏重魏畏魏。韩必不敢反魏。是韩则魏之县也。魏得韩以为县。卫。大梁。河外必安矣。今不存韩。二周。安陵必危。楚。赵大破。卫。齐甚畏。天下西乡而驰秦入朝而为臣不久矣。\\
\hspace*{2em}二十年。秦围邯郸。信陵君无忌矫夺将军晋鄙兵以救赵。赵得全。无忌因留赵。二十六年。秦昭王卒。\\
\hspace*{2em}三十年。无忌归魏。率五国兵攻秦。败之河外。走蒙骜。魏太子增质於秦。秦怒。欲囚魏太子增。或为增谓秦王曰。﹁公孙喜固谓魏相曰﹃请以魏疾击秦。秦王怒。必囚增。魏王又怒。击秦。秦必伤﹄。今王囚增。是喜之计中也。故不若贵增而合魏。以疑之於齐。韩。﹂秦乃止增。\\
\hspace*{2em}三十一年。秦王政初立。\\
\hspace*{2em}三十四年。安釐王卒。太子增立。是为景湣王。信陵君无忌卒。\\
\hspace*{2em}景湣王元年。秦拔我二十城。以为秦东郡。二年。秦拔我朝歌。徙野王。三年。秦拔我汲。五年。秦拔我垣。蒲阳。衍。十五年。景湣王卒。子王假立。\\
\hspace*{2em}王假元年。燕太子丹使荆轲刺秦王。秦王觉之。\\
\hspace*{2em}三年。秦灌大梁。虏王假。遂灭魏以为郡县。\\
\hspace*{2em}太史公曰。吾適故大梁之墟。墟中人曰。﹁秦之破梁。引河沟而灌大梁。三月城坏。王请降。遂灭魏。﹂说者皆曰魏以不用信陵君故。国削弱至於亡。余以为不然。天方令秦平海内。其业未成。魏虽得阿衡之佐。曷益乎。\\

% ===== 索隐述赞（司马贞《史记索隐》）· 统一样式 =====
% 题头「索隐述赞曰」五字并入反白框；正文每四句一列，定宽盒 \hfil 居列中；
% \fontsize{24}{41} 列距对齐全书界行（121+41k）；\clearpage 使述赞独占一页。
\clearpage
\shuzanhead{索隐述赞曰}\linebreak
{\shuzanbody\fontsize{24}{41}\selectfont
\hbox to \hsize{\hfil 毕公之苗。因国为姓。大名始赏。盈数自正。\hfil}\linebreak
\hbox to \hsize{\hfil 胤裔繁昌。系载忠正。杨干就戮。智氏奔命。\hfil}\linebreak
\hbox to \hsize{\hfil 文始建侯。武实彊盛。大梁东徙。长安北侦。\hfil}\linebreak
\hbox to \hsize{\hfil 卯既无功。卬亦外聘。王假削弱。虏於秦政。\hfil}%
\par}
\end{document}
